%% kapa: exam=1 round=30 year=2019 period=1 form=A subject=real_estate subject_ko=부동산학원론 seq=18 number=98
%% source: Q-Net 감정평가사 기출문제 (cst00309) artlSeq=5205205 file=extracted/1008/5205205/2019년 제30회 감정평가사 1차시험 문제지/2019년 제30회 감정평가사 1차시험 문제지/2019년도 제30회 감정평가사 1차 1교시 A형.hwp
%% conversion: kordoc 4.24.51 -> tools/gen1.py (auto); review=proofread-sonnet; flags=none
\begin{question}{98}
부동산조세에 관한 설명으로 옳지 않은 것은? (단, 주어진 조건에 한함)
\begin{choices}
  \item 종합부동산세와 재산세의 과세대상은 일치한다.
  \item 조세의 귀착 문제는 수요와 공급의 상대적 탄력성에 달려 있다.
  \item 임대주택에 재산세가 강화되면 장기적으로 임차인에게 전가될 수 있다.
  \item 부동산조세는 자원을 재분배하는 기능이 있다.
  \item 주택에 보유세가 중과되면 자가소유 수요가 감소할 수 있다.
\end{choices}
\end{question}
\answer{1}
